%%%%%%%% ICML 2026 WORKSHOP SUPPLEMENTARY MATERIAL %%%%%%%%%%%%%%%%%%%%%%%%
% Standalone supplementary appendix for the 4-page main paper.
% Submit alongside CausalCellBench_ICML_Short4pp.pdf as supplementary material.
\documentclass{article}

\usepackage{microtype}
\usepackage{graphicx}
\usepackage{subcaption}
\usepackage{booktabs}
\usepackage{multirow}
\usepackage{enumitem}
\usepackage{amsmath, amssymb, amsthm}
\usepackage{xcolor}
\usepackage{url}
\usepackage[colorlinks=true,
            linkcolor=blue,
            citecolor=blue,
            urlcolor=blue]{hyperref}

\usepackage{icml2026}

\usepackage{mathtools}
\usepackage[capitalize,noabbrev]{cleveref}

\icmltitlerunning{Virtual cell models -- Supplementary material}

\begin{document}

\twocolumn[
\icmltitle{Supplementary Material:\\
           Virtual Cell Models Inflate Perturbation Effect Sizes\\
           and Undermine Causal Gene Regulatory Network Recovery}

\icmlsetsymbol{equal}{*}

\begin{icmlauthorlist}
\icmlauthor{Anonymous Authors}{anon}
\end{icmlauthorlist}

\icmlaffiliation{anon}{Anonymous Institution}

\icmlkeywords{Virtual Cell Models, Causal Discovery, Perturb-seq, Gene Regulatory Networks, Foundation Models, Calibration}

\vskip 0.3in
]

\printAffiliationsAndNotice{}

\onecolumn
\appendix

\noindent\textbf{Notation and scope.} This supplementary material accompanies the 4-page main paper. Section and table references with prefix \texttt{tab:} or \texttt{fig:} refer to the main paper; section references with the prefix \texttt{app:} refer to the appendix below. All numbers are reproduced from the JSON artifacts in the released code repository, no values are paraphrased or rederived in prose.

\section{Additional results}

\subsection{Scale sweep precision across $N{=}50$--$200$}

\begin{table}[h]
\centering
\small
\begin{tabular}{rrrrrr}
\toprule
\textbf{$N$ (genes)} & \textbf{GT edges} & \textbf{ElasticNet} & \textbf{GEARS} & \textbf{CPA} & \textbf{Geneformer} \\
\midrule
50  & 169   & 0.345 & 0.280 & 0.302 & 0.292 \\
75  & 309   & 0.309 & 0.351 & 0.344 & 0.316 \\
100 & 509   & 0.344 & 0.339 & 0.335 & 0.348 \\
125 & 646   & 0.356 & 0.360 & 0.326 & 0.349 \\
150 & 821   & 0.368 & 0.310 & 0.347 & \textbf{0.394} \\
175 & 1{,}002 & 0.369 & \textbf{0.401} & 0.378 & 0.340 \\
200 & 1{,}217 & \textbf{0.378} & 0.392 & 0.385 & 0.373 \\
\bottomrule
\end{tabular}
\caption{Scale-sweep precision against real-data GIES ground truth at seven gene-set sizes. Bold values indicate top performance per row.}
\end{table}

\subsection{Multi-hop credit at $k = 1, 2, 3$}

\begin{table}[h]
\centering
\small
\begin{tabular}{lrrrrr}
\toprule
\textbf{Method} & \textbf{P@$k{=}1$} & \textbf{P@$k{=}2$} & \textbf{P@$k{=}3$} & \textbf{Boost (k=2)} & \textbf{Boost (k=3)} \\
\midrule
ElasticNet  & 0.425 & 0.540 & 0.543 & +11.5 & +11.9 \\
GEARS       & 0.410 & 0.529 & 0.533 & \textbf{+11.8} & +12.2 \\
CPA         & 0.387 & 0.500 & 0.508 & +11.3 & +12.0 \\
Geneformer  & 0.432 & 0.525 & 0.529 & +9.2  & +9.7  \\
Overlay Real & 0.423 & 0.533 & 0.539 & +11.0 & +11.6 \\
\bottomrule
\end{tabular}
\caption{Multi-hop credit precision at hop-$k$ over the $N{=}200$ K562 evaluation. Boost columns show percentage-point increase from $k{=}1$.}
\end{table}

\subsection{RPE1 cross-context structural comparison}

\begin{table}[h]
\centering
\small
\begin{tabular}{lrr}
\toprule
\textbf{Property} & \textbf{K562} & \textbf{RPE1} \\
\midrule
First singular vector variance     & 56.8\% & 34.5\% \\
Top-5 singular variance            & 69.8\% & 68.5\% \\
Frac near-zero effects             & 38.6\% & 15.4\% \\
Frac large effects ($|\log_2$FC$|>0.5$) & 2.7\%  & 17.3\% \\
GIES edges (real data)             & 1{,}220 & 810 \\
\bottomrule
\end{tabular}
\caption{Structural properties of real perturbation-response matrices for K562 and RPE1 at $N{=}200$.}
\end{table}

\subsection{Multi-seed bootstrap confidence}
\label{app:multi_seed}

Across five seeds (42, 123, 456, 789, 1024) the four predictive methods sit within overlapping intervals on every metric (Table~\ref{tab:multiseed}). The GRNBoost2 random-baseline floor confirms that all four predictive methods recover real signal substantially above chance, while none of the deep models significantly separates from ElasticNet at this seed budget.

\begin{table}[h]
\centering
\small
\begin{tabular}{lrrrr}
\toprule
\textbf{Method} & \textbf{F1} & \textbf{Precision} & \textbf{Recall} & \textbf{AUPRC} \\
\midrule
ElasticNet            & $0.426 \pm 0.016$ & $0.420 \pm 0.014$ & $0.431 \pm 0.019$ & $0.425 \pm 0.016$ \\
CPA                   & $0.412 \pm 0.016$ & $0.407 \pm 0.013$ & $0.417 \pm 0.019$ & $0.400 \pm 0.022$ \\
GEARS                 & $0.406 \pm 0.015$ & $0.404 \pm 0.012$ & $0.408 \pm 0.018$ & $0.419 \pm 0.012$ \\
Geneformer            & $0.425 \pm 0.011$ & $0.422 \pm 0.010$ & $0.429 \pm 0.013$ & $0.422 \pm 0.013$ \\
Overlay-real (oracle) & $0.422 \pm 0.009$ & $0.420 \pm 0.010$ & $0.425 \pm 0.009$ & $0.420 \pm 0.013$ \\
GRNBoost2 (random)    & $0.098 \pm 0.004$ & $0.074 \pm 0.003$ & $0.147 \pm 0.005$ & $0.077 \pm 0.005$ \\
\bottomrule
\end{tabular}
\caption{Multi-seed mean and standard deviation across $n=5$ seeds at $N{=}200$ on K562 against real-data GIES ground truth (1{,}220 edges per seed). Overlay-real is the same overlay-sampling pipeline applied to real perturbed cell predictions and serves as a noise-floor for the overlay step itself. GRNBoost2 is a structurally-distant baseline (gradient-boosted regression) included as a near-random reference.}
\label{tab:multiseed}
\end{table}

\subsection{Full PC and NOTEARS robustness tables}
\label{app:pc_notears}

The main-text PC and NOTEARS robustness paragraph cites summary statistics. Tables~\ref{tab:pc_full} and~\ref{tab:notears_full} report the full per-condition F1, precision, recall, and edge count against each backend's own real-data ground truth.

\begin{table}[h]
\centering
\small
\begin{tabular}{lrrrr}
\toprule
\textbf{Condition} & \textbf{Edges} & \textbf{F1} & \textbf{Precision} & \textbf{Recall} \\
\midrule
Real (PC ground truth)   & 682 & --- & --- & --- \\
GEARS-vanilla            & 556 & 0.0549 & 0.0612 & 0.0499 \\
GEARS-calibrated         & 652 & 0.0480 & 0.0491 & 0.0469 \\
CPA-vanilla              & 268 & 0.0800 & 0.1418 & 0.0557 \\
CPA-calibrated           & 294 & 0.0779 & 0.1293 & 0.0557 \\
Geneformer-vanilla       & 142 & 0.0097 & 0.0282 & 0.0059 \\
Geneformer-calibrated    & 240 & 0.0043 & 0.0083 & 0.0029 \\
\bottomrule
\end{tabular}
\caption{PC algorithm (Fisher-Z conditional independence at $\alpha{=}0.05$) per-condition results on K562 ACE matrices at $N{=}200$, seed 42.}
\label{tab:pc_full}
\end{table}

\begin{table}[h]
\centering
\small
\begin{tabular}{lrrrr}
\toprule
\textbf{Condition} & \textbf{Edges} & \textbf{F1} & \textbf{Precision} & \textbf{Recall} \\
\midrule
Real (NOTEARS ground truth) & 328 & --- & --- & --- \\
GEARS-vanilla            & 24  & 0.0511 & 0.3750 & 0.0274 \\
GEARS-calibrated         & 0   & 0.0000 & 0.0000 & 0.0000 \\
CPA-vanilla              & 18  & 0.0405 & 0.3889 & 0.0213 \\
CPA-calibrated           & 0   & 0.0000 & 0.0000 & 0.0000 \\
Geneformer-vanilla       & 7   & 0.0119 & 0.2857 & 0.0061 \\
Geneformer-calibrated    & 0   & 0.0000 & 0.0000 & 0.0000 \\
\bottomrule
\end{tabular}
\caption{NOTEARS (continuous-optimization with acyclicity constraint) per-condition results on K562 at $N{=}200$, seed 42, $\lambda{=}0.05$, weight threshold $w_{\mathrm{thresh}}{=}0.1$.}
\label{tab:notears_full}
\end{table}

\subsection{Six-variant CPA architectural ablation}
\label{app:cpa_ablation}

Six CPA variants were trained, each modifying one architectural axis relative to the NB-loss baseline. Per-variant causal F1 is reported in Table~\ref{tab:ablation_cpa}. Spread across all six variants is $0.387$--$0.418$ ($\Delta = 0.031$ F1), within the multi-seed CI of vanilla CPA.

\begin{table}[h]
\centering
\small
\begin{tabular}{lrrrr}
\toprule
\textbf{Variant} & \textbf{Edges} & \textbf{F1} & \textbf{Precision} & \textbf{Recall} \\
\midrule
\texttt{loss\_gauss}\,\,(Gaussian, not NB)        & 1{,}215 & 0.401 & 0.402 & 0.400 \\
\texttt{latent\_16}\,\,(latent dim 16, base 64)  & 1{,}188 & 0.413 & 0.418 & 0.407 \\
\texttt{latent\_256}\,\,(latent dim 256)         & 1{,}200 & 0.418 & 0.422 & 0.415 \\
\texttt{decoder\_1layer}\,\,(decoder depth 1)    & 1{,}207 & 0.418 & 0.420 & 0.416 \\
\texttt{layer\_norm}\,\,(LN replaces BN)         & 1{,}212 & 0.399 & 0.402 & 0.396 \\
\texttt{no\_adversarial}\,\,(reg=pen=0)          & 1{,}207 & 0.387 & 0.389 & 0.385 \\
\bottomrule
\end{tabular}
\caption{Six-variant CPA architectural ablation on K562 at $N{=}200$, seed 42, partition size 20, against real-data GIES ground truth (1{,}220 edges). All other hyperparameters are held at the vanilla CPA baseline (\texttt{n\_latent=64}, encoder layers $=3$, decoder layers $=3$, batch-norm, NB reconstruction loss, adversarial covariate disentanglement, max\_epochs $=300$).}
\label{tab:ablation_cpa}
\end{table}

\section{Supplementary methods}
\label{app:supp_methods}

\subsection{Virtual cell model hyperparameters}
\label{app:hparams}

\begin{table}[h]
\centering
\small
\begin{tabular}{ll}
\toprule
\textbf{Hyperparameter} & \textbf{Value} \\
\midrule
\multicolumn{2}{l}{\textit{CPA} (Lotfollahi et al., 2023)} \\
\hspace{1em}\texttt{n\_latent}                            & 64 \\
\hspace{1em}\texttt{n\_layers\_encoder}                   & 3 \\
\hspace{1em}\texttt{n\_layers\_decoder}                   & 3 \\
\hspace{1em}\texttt{batch\_size}                          & 512 \\
\hspace{1em}\texttt{max\_epochs}                          & 300 \\
\hspace{1em}\texttt{n\_epochs\_kl\_warmup}                & 20 \\
\hspace{1em}\texttt{n\_epochs\_adv\_warmup}               & 50 \\
\hspace{1em}\texttt{n\_epochs\_pretrain\_ae}              & 10 \\
\hspace{1em}reconstruction loss                           & Negative Binomial \\
\midrule
\multicolumn{2}{l}{\textit{GEARS} (Roohani et al., 2024)} \\
\hspace{1em}\texttt{hidden\_size}                          & 64 \\
\hspace{1em}\texttt{epochs}                                & 20 \\
\hspace{1em}\texttt{batch\_size}                           & 32 \\
\midrule
\multicolumn{2}{l}{\textit{Geneformer V2} (Theodoris et al., 2023)} \\
\hspace{1em}backbone                                       & \texttt{ctheodoris/Geneformer} (HF) \\
\hspace{1em}fine-tuning task                               & in-silico perturbation \\
\hspace{1em}\texttt{forward\_batch\_size}                  & 16 \\
\midrule
\multicolumn{2}{l}{\textit{ElasticNet baseline} (Zou \& Hastie, 2005)} \\
\hspace{1em}\texttt{alpha}                                 & 0.1 \\
\hspace{1em}\texttt{l1\_ratio}                             & 0.5 \\
\hspace{1em}\texttt{max\_iter}                             & 1{,}000 \\
\midrule
\multicolumn{2}{l}{\textit{STATE} (Arc Institute, 2025)} \\
\hspace{1em}embedding model                                & SE-600M \\
\hspace{1em}prediction model                               & ST-SE-Replogle / zeroshot/k562 \\
\hspace{1em}\texttt{hidden\_size}                          & 328 \\
\hspace{1em}\texttt{num\_attention\_heads}                 & 12 \\
\hspace{1em}\texttt{head\_dim}                             & 64 (decoupled, $\neq H/h$) \\
\hspace{1em}\texttt{num\_hidden\_layers}                   & 8 \\
\hspace{1em}\texttt{intermediate\_size}                    & 3{,}072 \\
\hspace{1em}\texttt{max\_position\_embeddings}             & 64 \\
\hspace{1em}embedding dim ($X_{\mathrm{state}}$)           & 2{,}058 \\
\bottomrule
\end{tabular}
\caption{Training hyperparameters per virtual cell model. Geneformer is fine-tuned from the public \texttt{ctheodoris/Geneformer} V2 checkpoint; STATE is run zero-shot from the released \texttt{ST-SE-Replogle} K562 checkpoint.}
\label{tab:hparams}
\end{table}

\subsection{STATE inference pipeline}
\label{app:state_methods}

The STATE result was obtained zero-shot from the released ST-SE-Replogle/zeroshot/k562 checkpoint via the \texttt{state} CLI v0.10.5. The pipeline has two stages: (1)~SE-600M cell encoder (\texttt{state emb transform}) maps each input cell to a 2{,}058-dimensional $X_{\mathrm{state}}$ embedding; (2)~the ST-SE-Replogle perturbation model (\texttt{state tx infer}) maps $(X_{\mathrm{state}}, \mathrm{perturbation})$ pairs to predicted gene-expression vectors over the same $N{=}200$ gene panel.

Reproducing the published checkpoint required diagnosing three distinct framework-versioning bugs:
\begin{enumerate}[leftmargin=1.5em,topsep=2pt,itemsep=2pt]
    \item \textbf{Decoupled head dimension.} The released checkpoint encodes a Llama backbone with \texttt{hidden\_size=328}, \texttt{num\_attention\_heads=12}, and \emph{explicit} \texttt{head\_dim=64}, i.e.\ $328 \neq 12 \times 64 = 768$. Transformers $\geq$5 add a strict \texttt{validate\_architecture} check that rejects \texttt{hidden\_size \% num\_attention\_heads != 0} regardless of an explicit \texttt{head\_dim}; transformers 4.45 accepts decoupled head dimensions but lacks the (2) compatibility we describe next; transformers 4.52.3 supports both. \textbf{Fix:} pin \texttt{transformers==4.52.3} after \texttt{pip install -e state}.
    \item \textbf{is\_causal\,kwarg pass-through.} The state package's \texttt{LlamaBidirectionalModel.forward} sets \texttt{is\_causal=False} and forwards it to the parent \texttt{LlamaModel.forward} via \texttt{**flash\_attn\_kwargs}. transformers 4.45's signature lacks the \texttt{**flash\_attn\_kwargs} catch-all, so the call raises \texttt{TypeError: unexpected keyword argument 'is\_causal'}. transformers 4.52.3 introduces the catch-all and accepts the call.
    \item \textbf{ALL\_PARALLEL\_STYLES\,is None on Modal.} In transformers 4.52.3 specifically, \texttt{transformers.modeling\_utils.ALL\_PARALLEL\_STYLES} is initialized to \texttt{None} in some container configurations (verified locally and on Modal), making \texttt{LlamaModel.\_\_init\_\_}'s \texttt{post\_init} hit \texttt{if v not in ALL\_PARALLEL\_STYLES} and raise \texttt{TypeError: argument of type 'NoneType' is not iterable}. \textbf{Fix:} during image build, replace the offending line in \texttt{transformers/modeling\_utils.py} with \texttt{if ALL\_PARALLEL\_STYLES is not None and v not in ALL\_PARALLEL\_STYLES:}; this propagates to the \texttt{state} CLI subprocess, where a \texttt{sitecustomize.py} monkey-patch did not reliably fire.
\end{enumerate}

The full set of attempted patches and the final working Modal image specification are in \texttt{scripts/modal\_run\_state\_v3.py} of the released code.

\textbf{Subset rationale.} STATE inference was run on a 320-cell subset of K562: $120$ control cells (sampled from the $\sim$10{,}691 \texttt{non-targeting} cells) plus $200$ perturbed cells (one per gene in the $N{=}200$ panel). The reduction is required because SE-600M wall-clock inference on Modal CPU instances measures $\sim$28\,s/cell at \texttt{batch\_size=8}; the full 38{,}979-cell K562 panel would take $\sim$300\,h, which exceeds the 6-hour Modal function timeout. The 320-cell subset embeds in $\sim$13.5 minutes and produces 200 distinct STATE-predicted perturbation profiles---one per perturbation, the unit consumed by overlay sampling and downstream GIES.

\subsection{Norman 2019 cross-paradigm pipeline}
\label{app:norman_pipeline}

The Norman 2019 dataset (GSE133344) was downloaded from CellxGene Census and contains $111{,}445$ K562 cells under $237$ unique perturbations. After dropping combinatorial perturbations and keeping only \texttt{control} (n=11{,}855) plus single-gene perturbations with at least 10 cells, $102$ eligible target genes remain. The substitutability evaluation re-uses the K562 protocol unchanged: 200 control cells + 100 cells per perturbation overlay-sampled from the ElasticNet predictions, vstacked into one expression matrix and fed to GIES with partition size 20 and seed 42.

\textbf{Two practical engineering issues are noted.} (i)~\texttt{multiprocessing.Pool(fork)} for parallel partition execution deadlocks on the semaphore in Modal's container runtime; switching to sequential GIES execution removes the deadlock at the cost of single-thread CPU throughput (acceptable: total wall time ${\approx}28$\,s on the M4 laptop). (ii)~Some Norman partitions contain genes with zero variance under particular interventional regimes, producing \texttt{1/omega = inf} in the gauss\_int score and an effectively infinite hill-climbing loop in GIES. Adding $\mathcal{N}(0, 10^{-6})$ jitter to the partition expression matrix before scoring prevents the degeneracy without measurably altering recovered edge sets in the well-conditioned partitions.

\subsection{Causal-discovery backend implementations}
\label{app:cd_backends}

\textbf{GIES}: Greedy Interventional Equivalence Search via the \texttt{gies} Python package (v0.0.1), score function \texttt{gauss\_int\_l0\_pen}, ${\sim}20$-gene partitions for tractability of the underlying CPDAG search. Output edges are the union of single-direction CPDAG arrows across partitions; bidirectional CPDAG arrows are kept once with the canonical ordering (lower index first).

\textbf{inspre}: Sparse-inverse regression on the perturbation-conditional ACE matrix via the \texttt{inspre} R package, full $\lambda$-path with leave-one-target-out cross-validation. The Modal R kernel (renv-isolated) is invoked via \texttt{Rscript} subprocess to keep the Python pipeline driver-only.

\textbf{PC}: Constraint-based causal discovery via the \texttt{causal-learn} Python package, Fisher-Z conditional independence at $\alpha = 0.05$. Run on each model's per-perturbation ACE matrix.

\textbf{NOTEARS}: Continuous-optimization with acyclicity constraint, $\lambda = 0.05$, weight threshold $w_{\mathrm{thresh}} = 0.1$.

\textbf{Overlay sampling.} For each perturbed condition, $100$ synthetic cells are produced by drawing $100$ control cells with replacement, multiplying by a per-gene ratio of (predicted-mean $+ 10^{-2}$) over (control-mean $+ 10^{-2}$) clipped to $[0.01, 10]$, and Poisson-resampling to count-like expression. This preserves per-cell control covariance while imposing the predicted shift in marginal magnitude.

\section{Reproducibility}
\label{app:reproducibility}

\subsection{Software environment}
\label{app:software}

\begin{table}[h]
\centering
\small
\begin{tabular}{ll}
\toprule
\textbf{Component} & \textbf{Version} \\
\midrule
Python                                & 3.9 (local) / 3.11 (Modal) \\
PyTorch                               & 2.4.1 (CPU) \\
\texttt{transformers}                 & 4.52.3 (STATE-pinned) \\
\texttt{tokenizers}                   & 0.21.x \\
\texttt{huggingface\_hub}             & 0.30.x \\
\texttt{anndata}                      & 0.10.9 \\
\texttt{scanpy}                       & 1.10.3 \\
\texttt{numpy}                        & 1.26.4 \\
\texttt{scipy}                        & 1.13.1 \\
\texttt{scikit-learn}                 & 1.x \\
\texttt{gies} (Python)                & 0.0.1 \\
\texttt{causal-learn}                 & 0.1.x \\
\texttt{notears}                      & reference impl. \\
\texttt{inspre} (R)                   & 0.0.1 \\
\texttt{state} (Arc Institute)        & 0.10.5 \\
\bottomrule
\end{tabular}
\caption{Software versions used end-to-end. Modal containers pin all dependencies via the image build steps in \texttt{scripts/modal\_run\_*.py}.}
\label{tab:software}
\end{table}

\subsection{Compute resources}
\label{app:compute}

Local pipeline (overlay sampling, GIES, PC, NOTEARS, calibration, multi-seed aggregation, plotting) runs on an Apple M4 laptop with 16\,GB unified memory. Wall time: $\approx$45 minutes for the full $N{=}200$, single-seed K562 evaluation.

Modal cloud is used for the model-training and large-scale inference steps: GEARS, CPA, Geneformer, STATE, and the inspre $N{=}500$ scale point. Modal H100 instances are reserved for GEARS and Geneformer fine-tuning ($\approx$1\,h each at $N{=}200$); A10G for CPA training ($\approx$45 minutes); 16-CPU 64\,GB instances for inspre and the STATE CPU fallback. The largest single Modal job wall-clock is the SE-600M embedding step in the STATE pipeline at $\approx$3.5 hours for the full embedding, $\approx$13 minutes for the 320-cell subset.

\subsection{Random seeds}
\label{app:seeds}

The single-seed numbers throughout the main text use \texttt{seed=42}. Multi-seed CI tables in \S\ref{app:multi_seed} use $\{42, 123, 456, 789, 1024\}$. Inside the GIES partitioning step, the gene-permutation seed is $42$; partition contents are deterministic given the panel and seed.

\subsection{End-to-end reproduction commands}
\label{app:commands}

All numbers in the main text and tables are reproducible from the released code with \texttt{seed=42}.

\begin{verbatim}
# Calibrate predictions
PYTHONPATH=. python3 scripts/calibrate_predictions.py \
    --diagnostics-out results/calibration_diagnostics.json

# Vanilla evaluation
PYTHONPATH=. OMP_NUM_THREADS=1 python3 \
  scripts/run_eval_local.py \
    --data-path data/k562_subset_n200_slim.h5ad \
    --pre-subset --n-genes 200 --partition-size 20 \
    --n-cells 100 --n-ctrl 200 --seed 42 \
    --model-outputs-dir data/model_outputs_real_n200 \
    --gies-cache-dir results/gies_cache --skip-random

# Calibrated evaluation
PYTHONPATH=. OMP_NUM_THREADS=1 python3 \
  scripts/run_eval_local.py \
    --data-path data/k562_subset_n200_slim.h5ad \
    --pre-subset --n-genes 200 --partition-size 20 \
    --n-cells 100 --n-ctrl 200 --seed 42 \
    --model-outputs-dir \
       data/model_outputs_real_n200_calibrated \
    --gies-cache-dir results/gies_cache_calibrated \
    --skip-random

# STATE evaluation
PYTHONPATH=. OMP_NUM_THREADS=1 python3 \
  scripts/state_eval_gies.py

# Norman 2019 cross-paradigm
.venv_gears/bin/python scripts/run_norman_local.py
\end{verbatim}

\end{document}
